% Einheit 3 von 4 – Multiplizieren und Dividieren (bank/rationale-zahlen/e3.jsonl, zusammenbau v0.1)
%% TODO Zweigzeile Teil 1: was hier gelernt wird (Satz in Schülersprache aus dem Einheitstitel) – Einheit 3
\einheitenkopf[e3]{Einheit 3 von 4 · Multiplizieren und Dividieren}
\zweigzeile{ab Kl. 7 · P10}
\setcounter{aufgabe}{19}

%% TODO Titel: Ich-kann-Satz fehlt in der Bank (Platzhalter: Kettenname) – Nr. 20
%% TODO Anweisung: ein Satz über der ersten Teilaufgabe fehlt in der Bank – Nr. 20
\begin{aufgabe}{Multiplizieren/Dividieren}
\begin{teile}
\teil $(-6) \cdot 7$ – Vorzeichen des Ergebnisses?\\ \kreuz{plus} \kreuz{minus}
\end{teile}
%% TODO Darstellung für schwach fehlt in der Bank (rationale-zahlen-e3-k1-s1-v1; Abschnitt „Für schwache Schüler“)
\swz{Berechne: $7 \cdot (-5)$}{}
%% TODO Darstellung für schwach fehlt in der Bank (rationale-zahlen-e3-k1-s1-v2; Abschnitt „Für schwache Schüler“)
\swz{Berechne: $(-8) \cdot 6$}{}
%% TODO Darstellung für schwach fehlt in der Bank (rationale-zahlen-e3-k1-s1-v3; Abschnitt „Für schwache Schüler“)
\swz{Berechne: $9 \cdot (-3)$}{}
%% TODO Darstellung für schwach fehlt in der Bank (rationale-zahlen-e3-k1-s1-v4; Abschnitt „Für schwache Schüler“)
\swz{Berechne: $(-15) \cdot 4$}{}
%% TODO Darstellung für schwach fehlt in der Bank (rationale-zahlen-e3-k1-s1-v5; Abschnitt „Für schwache Schüler“)
\swz{Der Wasserstand sinkt jeden Tag um $3$ cm. Wie ändert er sich in $6$ Tagen?}{}
\swfrage{Was bleibt gleich, was ändert sich?}
%% TODO Darstellung für schwach fehlt in der Bank (rationale-zahlen-e3-k1-s2-v1; Abschnitt „Für schwache Schüler“)
\swz{Berechne: $(-7) \cdot (-6)$}{}
%% TODO Darstellung für schwach fehlt in der Bank (rationale-zahlen-e3-k1-s3-v1; Abschnitt „Für schwache Schüler“)
\swz{Berechne: $(-56) : 8$}{}
%% TODO Darstellung für schwach fehlt in der Bank (rationale-zahlen-e3-k1-s4-v1; Abschnitt „Für schwache Schüler“)
\swz{Berechne: $(-2) \cdot 5 \cdot (-3)$}{}
%% TODO Darstellung für schwach fehlt in der Bank (rationale-zahlen-e3-k1-s5-v1; Abschnitt „Für schwache Schüler“)
\swz{Berechne: $(-4)^2$ und $-4^2$}{}
%% TODO Darstellung für schwach fehlt in der Bank (rationale-zahlen-e3-k1-s6-v1; Abschnitt „Für schwache Schüler“)
\swz{Berechne: $5 + 3 \cdot (-4)$}{}
\swz[3]{Berechne ohne Taschenrechner den Wert von $\frac{a + b}{c}$ für $a = 15$, $b = -6$ und $c = -4$. \hfill (P10 2021 OS)}{}
\swz[3]{Berechne den Wert von $(a + b) : c$ für $a = 3$, $b = -9$ und $c = -2$. \hfill (P10 2016 OS)}{}
\swz[3]{Berechne den Wert von $3 \cdot (x - 5)$ für $x = -4$. \hfill (P10 2026 FOR)}{}
\begin{teile}
\teil Eine negative Zahl wird mit einer negativen Zahl multipliziert. Das Ergebnis ist … \hfill (P10 2015 OS)\\ \kreuz{immer positiv}\\ \kreuz{immer negativ}\\ \kreuz{Das kann man nicht entscheiden.}
\end{teile}
\end{aufgabe}

%% TODO Titel: Ich-kann-Satz fehlt in der Bank (Platzhalter: Kettenname) – Nr. 21
%% TODO Anweisung: ein Satz über der ersten Teilaufgabe fehlt in der Bank – Nr. 21
\begin{aufgabe}{Vorzeichenregel als Aussage prüfen}
\begin{teile}
\teil Stimmt es? Das Produkt aus drei negativen Zahlen ist positiv. \janein
\end{teile}
\end{aufgabe}

%% TODO Titel: Ich-kann-Satz fehlt in der Bank (Platzhalter: Kettenname) – Nr. 22
%% TODO Anweisung: ein Satz über der ersten Teilaufgabe fehlt in der Bank – Nr. 22
\begin{aufgabe}{Multiplizieren/Dividieren – Fehler finden, Begründen, Anwendung}
\swz[3]{Nina rechnet: $(-6) \cdot 5 = 30$. Finde den Fehler.}{}
\swfrage{Setze fort: $2 \cdot (-4) = -8$, $1 \cdot (-4) = -4$, $0 \cdot (-4) = 0$. Begründe damit, warum $(-1) \cdot (-4)$ positiv ist.}
\swz[3]{Ein Handyvertrag bucht jeden Monat $15$ € ab. Wie ändert sich der Kontostand in einem halben Jahr?}{}
\end{aufgabe}

\uebersichtskasten{Vorzeichenregel: gleiche Vorzeichen → plus, verschiedene → minus. Gilt für mal und geteilt. \\ (−3) · 8 = −24 \quad (−3) · (−8) = 24 \quad 24 : (−8) = −3 \\ Mehrere Faktoren: gerade Anzahl Minus → plus, ungerade → minus. \quad (−1) · (−1) · (−1) = −1 \\ Potenz: (−5)² = 25, aber −5² = −25 (das Minus gehört nicht zur Basis).}
